% =========================================================================
% SECCIÓN 3: ANÁLISIS DE EJES TEMÁTICOS Y TECNOLÓGICOS
% =========================================================================

\section{Análisis de Ejes Temáticos y Tecnológicos}
\label{sec:ejes_tematicos}

\subsection{Eje 1: Modelos Compactos de Lenguaje (SLMs) y Despliegue Local}

\subsubsection{Familias de modelos compactos de última generación}
\noindent [Análisis comparativo de modelos base ligeros: Qwen2.5 (0.5B, 1.5B, 3B), LLaMA-3.2 (1B, 3B), Phi-3.5-mini, Gemma-2-2B.]

\subsubsection{Técnicas de compresión y cuantización}
\noindent [Revisión de métodos de cuantización post-entrenamiento: GGUF/GGML (Q4\_K\_M, Q5\_K\_M), AWQ y GPTQ; impacto en perplejidad y uso de VRAM.]

\subsubsection{Motores de inferencia offline en hardware restringido}
\noindent [Desempeño de frameworks como llama.cpp, Ollama, ONNX Runtime y vLLM en CPUs estándar y GPUs con 4\,GB de VRAM.]

\subsection{Eje 2: Generación Aumentada por Recuperación (RAG)}

\subsubsection{Estrategias de fragmentación y chunking contextual}
\noindent [Evaluación de segmentación semántica vs. por longitud de caracteres para preservar coherencia textual y morfológica.]

\subsubsection{Modelos de incrustación multilingüe e índices vectoriales}
\noindent [Análisis de modelos de embeddings locales (bge-m3, multilingual-e5, MiniLM) y gestores vectoriales embebidos (ChromaDB, FAISS, Qdrant).]

\subsubsection{Técnicas de reordenamiento (Re-ranking) y filtrado}
\noindent [Impacto de re-rankers compactos (bge-reranker) en la precisión del contexto recuperado para mitigar alucinaciones.]

\subsection{Eje 3: Procesamiento de Lenguaje Natural en Quechua Collao}

\subsubsection{Particularidades morfológicas y desafío de tokenización}
\noindent [Análisis de la naturaleza aglutinante y polisintética del quechua; subtokenización excesiva (token fertility) en tokenizadores BPE estándar.]

\subsubsection{Recursos lingüísticos y corpus bilingües disponibles}
\noindent [Inventario crítico de corpus texto-texto (Tatoeba, JW300, AmericasNLP), vocabularios normalizados y diccionarios oficiales MINEDU.]

\subsubsection{Translenguaje y alternancia de códigos (Code-switching)}
\noindent [Manejo de préstamos léxicos del castellano en el habla coloquial quechua andina en sistemas conversacionales.]

\subsection{Eje 4: Asistentes Educativos Basados en IA y Andamiaje Pedagógico}

\subsubsection{Sistemas conversacionales para Educación Intercultural Bilingüe (EIB)}
\noindent [Revisión de asistentes de tutoría inteligente orientados a la revitalización lingüística y mediación cultural en aulas rurales.]

\subsubsection{Diseño de prompts pedagógicos y andamiaje cognitivo}
\noindent [Estrategias de prompting (Socrático, Few-shot) para guiar al estudiante sin entregar respuestas directas.]

\subsection{Eje 5: Adaptación Eficiente en Parámetros (PEFT) y Alineación}

\subsubsection{Afinamiento con LoRA y QLoRA en hardware limitado}
\noindent [Mecanismos de adaptación de bajo rango sobre modelos cuantizados a 4 bits; presupuestos de memoria VRAM ($<$ 4\,GB).]

\subsubsection{Alineación de instrucciones y prevención de olvido catastrófico}
\noindent [Metodologías de instruction-tuning con datos sintéticos y reales; estrategias de regularización multilingüe.]
